\section{Finite-state framework and the two logarithms}\label{sec:framework}

\subsection{One-sided edge shifts and twisted determinants}

Throughout, \(\Gamma=(V,E,o,t)\) is a finite directed multigraph.  Parallel
edges are allowed and are regarded as named edges.  Its adjacency matrix is
\[
  A_{ij}=\#\{e\in E:o(e)=i,\ t(e)=j\}.
\]
We assume that \(A\) is primitive and that its Perron root satisfies
\(\lambda>1\).  The phase space in every main statement is the one-sided
edge shift
\[
  X_\Gamma^+=\{(e_0,e_1,\ldots):t(e_j)=o(e_{j+1})\},
  \qquad \sigma(e_0e_1\ldots)=e_1e_2\ldots.
\]
No assertion about arbitrary two-sided cocycles is implicit in this
notation.

Fix a finite group \(G\) and an edge labelling \(\tau:E\to G\).  For a
directed path \(q=e_0\cdots e_{n-1}\), set
\(\tau(q)=\tau(e_0)\cdots\tau(e_{n-1})\).  Changing the initial point on a
closed path conjugates this product.  Hence a periodic orbit \(\gamma\) has a
well-defined Frobenius conjugacy class \([g_\gamma]\), although a representative
\(g_\gamma\) requires a base point.

For a finite-dimensional unitary representation
\(\rho:G\to U(V_\rho)\), define the twisted adjacency matrix by blocks
\begin{equation}\label{eq:twisted-adjacency}
  (A_\rho)_{ij}:=\sum_{e:i\to j}\rho(\tau(e))
  \quad\in\End(V_\rho),
  \qquad
  P_\rho(z):=\det(I-zA_\rho).
\end{equation}
For the trivial representation, \(A_{\mathbf1}=A\).  Reversing every
representation convention replaces labels by inverses and classes by inverse
classes; none of the distinction between \(F\) and \(L\) depends on that
choice.

Let \(N_{n,C}\) be the number of period-\(n\) points whose length-\(n\)
holonomy lies in the conjugacy class \(C\).  Let \(p_{n,C}\) be the number of
primitive periodic orbits of least period \(n\) with Frobenius class \(C\).
For a class function \(\phi:G\to\CC\), use the abbreviations
\begin{align}
  T_n(\phi)&:=\sum_C\phi(C)N_{n,C},
  &c_n(\phi)&:=\sum_C\phi(C)p_{n,C},\label{eq:weighted-counts}\\
  \Pi_\phi(z)&:=\sum_{n\geq1}c_n(\phi)z^n
  =\sum_{\gamma\in\cP}\phi(g_\gamma)z^{|\gamma|}.
  \label{eq:primitive-channel}
\end{align}

\begin{lemma}[Twisted trace and periodic logarithm]\label{lem:trace-log}
For every irreducible representation \(\rho\) and every \(n\geq1\),
\begin{equation}\label{eq:twisted-trace}
  \Tr(A_\rho^n)=T_n(\chi_\rho).
\end{equation}
Consequently, on \(|z|<\lambda^{-1}\),
\begin{equation}\label{eq:artin-determinant-log}
  L_{\chi_\rho}(z):=-\Log_0P_\rho(z)
  =\sum_{n\geq1}\frac{T_n(\chi_\rho)}{n}z^n
  =\sum_{r\geq1}\frac{\Pi_{\psi^r\chi_\rho}(z^r)}{r}.
\end{equation}
Here \(\Log_0P_\rho\) is the analytic germ at zero normalized by
\(\Log_0P_\rho(0)=0\).
\end{lemma}

\begin{proof}
Expanding the block product in \(A_\rho^n\) sums \(\rho(\tau(q))\) over
closed length-\(n\) paths, with a trace over \(V_\rho\); this is
\eqref{eq:twisted-trace}.  The finite-matrix identity
\[
  -\Log_0\det(I-zB)=\sum_{n\geq1}\Tr(B^n)z^n/n
\]
gives the second expression in \eqref{eq:artin-determinant-log}.  A primitive
orbit \(\gamma\) of length \(d\) contributes its \(d\) cyclic base points to
period \(n=rd\), and its repeated holonomy is \(g_\gamma^r\).  Division by
\(n=rd\) leaves the factor \(1/r\), proving the last expression.  Absolute
convergence follows from the finite-matrix trace bound on every smaller disc.
\end{proof}

For a general class function, write
\[
  \phi=\sum_{\rho\in\Irr(G)}
     \widehat\phi(\rho)\chi_\rho,
  \qquad
  \widehat\phi(\rho)
    =\frac1{|G|}\sum_{g\in G}\phi(g)\overline{\chi_\rho(g)},
\]
and define \(L_\phi=\sum_\rho\widehat\phi(\rho)L_{\chi_\rho}\).  This agrees
with the orbit expansion in \eqref{eq:artin-determinant-log}.  The Adams
operation \(\psi^m\) acts on class functions by
\((\psi^m\phi)(g)=\phi(g^m)\).  It need not preserve irreducibility or
positivity, but it has a finite character expansion
\begin{equation}\label{eq:adams-coefficients}
  \psi^m\chi_\rho=\sum_{\pi\in\Irr(G)}
      a_{\rho,\pi}^{(m)}\chi_\pi,
  \qquad
  a_{\rho,\pi}^{(m)}
     =\langle\psi^m\chi_\rho,\chi_\pi\rangle_G.
\end{equation}
The coefficients are periodic in \(m\) modulo the exponent of \(G\), and
therefore uniformly bounded.

\subsection{Adams--M\"obius primitive inversion}

\begin{theorem}[Primitive and fixed-label determinant coordinates]
\label{thm:adams-mobius}
For every class function \(\phi\) and \(|z|<\lambda^{-1}\),
\begin{align}
  \Pi_\phi(z)
    &=\sum_{m\geq1}\frac{\mu(m)}{m}
       L_{\psi^m\phi}(z^m),
       \label{eq:primitive-adams-mobius}\\
  F_\phi(z):=\sum_{r\geq1}\frac{\Pi_\phi(z^r)}{r}
    &=\sum_{r,m\geq1}\frac{\mu(m)}{rm}
       L_{\psi^m\phi}(z^{rm}).
       \label{eq:fixed-adams-mobius}
\end{align}
For \(\phi=\chi_\rho\), insertion of \eqref{eq:adams-coefficients} gives a
formula involving only the finite determinant family:
\begin{equation}\label{eq:fixed-determinant-open-disc}
  F_{\chi_\rho}(z)=
  \sum_{r,m\geq1}\frac{\mu(m)}{rm}
    \sum_{\pi\in\Irr(G)}a_{\rho,\pi}^{(m)}
      \bigl(-\Log_0P_\pi(z^{rm})\bigr).
\end{equation}
All three series are absolutely convergent on compact subsets of the Perron
disc after expansion into primitive orbits.
\end{theorem}

\begin{proof}
By \Cref{lem:trace-log},
\[
  L_{\psi^m\phi}(z^m)
  =\sum_{\gamma\in\cP}\sum_{r\geq1}
     \frac{\phi(g_\gamma^{mr})}{r}z^{mr|\gamma|}.
\]
Multiplying by \(\mu(m)/m\), summing, and putting \(s=mr\), the coefficient
of \(\phi(g_\gamma^s)z^{s|\gamma|}\) becomes
\[
  \sum_{m\mid s}\frac{\mu(m)}{m(s/m)}
  =\frac1s\sum_{m\mid s}\mu(m).
\]
It is one for \(s=1\) and zero otherwise.  This proves
\eqref{eq:primitive-adams-mobius}.  Summing that identity over \(z^r/r\)
proves \eqref{eq:fixed-adams-mobius}.  Finally use the character expansion
\eqref{eq:adams-coefficients} and \eqref{eq:artin-determinant-log}.

For convergence, the total number \(p_n\) of primitive length-\(n\) orbits
satisfies \(p_n\leq\Tr(A^n)/n\ll\lambda^n/n\).  Thus the orbit expansions
are dominated on \(|z|\leq R<\lambda^{-1}\) by convergent logarithmic series.
This also justifies every rearrangement above.
\end{proof}

The difference between the periodic and fixed-label transforms is now
formal rather than terminological.

\begin{proposition}[Exact Artin-substitution obstruction]
\label{prop:artin-substitution-obstruction}
On \(|z|<\lambda^{-1}\),
\begin{equation}\label{eq:exact-f-minus-l}
  F_\phi(z)-L_\phi(z)
  =\sum_{r\geq2}\frac{
       \Pi_\phi(z^r)-\Pi_{\psi^r\phi}(z^r)}{r}.
\end{equation}
Writing \(\Pi_\phi(z)=\sum_{n\geq1}c_n(\phi)z^n\), the two germs agree if
and only if
\begin{equation}\label{eq:coefficient-obstruction}
  \sum_{r\mid n}\frac1r
    c_{n/r}(\phi-\psi^r\phi)=0
  \qquad(n\geq1).
\end{equation}
There is no general character-theoretic reason for these identities to hold.
\end{proposition}

\begin{proof}
Subtract the two orbit expansions in
\eqref{eq:intro-artin}--\eqref{eq:intro-fixed}; the \(r=1\) terms cancel.
Extracting the coefficient of \(z^n\) gives
\eqref{eq:coefficient-obstruction}.
\end{proof}

\subsection{The Perron boundary is derived from the determinant}

Put
\begin{equation}\label{eq:strict-gap}
  \eta:=\max_{\rho\in\Irr(G),\ \rho\neq\mathbf1}\rad(A_\rho).
\end{equation}
The main product theorem assumes \(\eta<\lambda\).  Under this hypothesis,
the \(rm=1\) terms in \eqref{eq:fixed-determinant-open-disc} are regular for
non-trivial \(\pi\), and every term with \(rm\geq2\) is an interior value
because
\begin{equation}\label{eq:interior-radius}
  \rad(\lambda^{-rm}A_\pi)\leq\lambda^{1-rm}<1.
\end{equation}

We record the boundary convention in a form that makes clear that no branch
or peripheral factor is an additional input coordinate.

\begin{proposition}[Canonical radial boundary construction]
\label{prop:derived-boundary}
Let \(P_\pi(z)=\det(I-zA_\pi)\).  The normalized germ
\(-\Log_0P_\pi(z)\), its continuation along \(z=t/\lambda\) with
\(0\leq t<1\), its peripheral factorization, and every radial renormalized
endpoint are determined by \(P_\pi\), \(\lambda\), and this path convention.
More explicitly, factor
\begin{equation}\label{eq:peripheral-factorization}
  P_\pi(t/\lambda)
   =\prod_{\zeta\in\Omega_\pi}(1-t\zeta)^{m_{\pi,\zeta}}
      Q_{\partial,\pi}(t),
\end{equation}
where \(\Omega_\pi\subset\TT\) is the multiset of peripheral phases and
\(Q_{\partial,\pi}\) has no zero on \([0,1]\).  If
\(s_\pi=m_{\pi,1}\), then
\begin{align}
  B_{1,\pi}^{\mathrm{ren}}
  &:={\lim}_{t\uparrow1}
     \left(-\Log_{[0,t]}P_\pi(t/\lambda)+s_\pi\log(1-t)\right)
     \label{eq:derived-b-ren}\\
  &=-\sum_{\substack{\zeta\in\Omega_\pi\\\zeta\neq1}}
       m_{\pi,\zeta}\Log_{[0,1]}(1-\zeta)
     -\Log_{[0,1]}Q_{\partial,\pi}(1).
     \label{eq:derived-b-factor}
\end{align}
In particular, a list of branch rules, peripheral factors, or values such as
\(B_{k,\pi}^{\mathrm{ren}}\) is derived output, not independent cocycle data.
\end{proposition}

\begin{proof}
The roots and their multiplicities in \eqref{eq:peripheral-factorization}
are determined by the polynomial \(P_\pi(t/\lambda)\).  The normalized
logarithm has a unique continuation along any zero-free subinterval beginning
at zero.  Removing the factor \((1-t)^{s_\pi}\), taking the radial limit, and
using the same continuation on every remaining factor yields
\eqref{eq:derived-b-factor}.  The construction for \(k>1\) is the identical
one applied at \(t^k/\lambda^k\); by \eqref{eq:interior-radius} it is an
ordinary interior logarithm in the strict-gap theorem below.
\end{proof}

\begin{lemma}[Perron evaluation of a non-trivial fixed-label coordinate]
\label{lem:perron-fixed-coordinate}
Assume \(\eta<\lambda\).  For every non-trivial irreducible \(\rho\), the
series
\begin{equation}\label{eq:perron-fixed-determinant}
  \mathfrak F_\rho:=
  \sum_{r,m\geq1}\frac{\mu(m)}{rm}
    \sum_{\substack{\pi\in\Irr(G)\\(r,m,\pi)\neq(1,1,\mathbf1)}}
       a_{\rho,\pi}^{(m)}
       \bigl(-\Log_{[0,1]}P_\pi(\lambda^{-rm})\bigr)
\end{equation}
is well defined and equals \(F_{\chi_\rho}(\lambda^{-1})\).  In
\eqref{eq:perron-fixed-determinant}, the notation means radial continuation
for the finitely many terms with \(rm=1\) and the normalized interior germ
for \(rm\geq2\).  The omitted scalar Perron triple has coefficient
\(a_{\rho,\mathbf1}^{(1)}=0\) and is deleted before endpoint evaluation.
\end{lemma}

\begin{proof}
Since \(\rho\neq\mathbf1\), character orthogonality gives
\(a_{\rho,\mathbf1}^{(1)}=0\).  If \(rm=1\), all remaining \(\pi\) are
non-trivial and \(P_\pi(t/\lambda)\neq0\) for \(0\leq t\leq1\) by the strict
gap.  For \(rm\geq2\), the trace expansion and
\(\abs{\Tr(A_\pi^n)}\leq(\dim A_\pi)\lambda^n\) give, uniformly for
\(0\leq t\leq1\),
\[
 \abs{-\Log_0P_\pi(t^{rm}\lambda^{-rm})}
 \leq C\sum_{n\geq1}\frac{\lambda^{(1-rm)n}}n
 \leq C'\lambda^{1-rm}.
\]
The Adams coefficients are uniformly bounded, and
\[
  \sum_{rm\geq2}\frac{\lambda^{1-rm}}{rm}
  =\sum_{k\geq2}\frac{d(k)}k\lambda^{1-k}<\infty.
\]
Dominated convergence in \eqref{eq:fixed-determinant-open-disc}, along
\(z=t/\lambda\), proves the assertion.
\end{proof}
